Domain fine-tuning erodes the general abilities of pretrained language models, and benchmarking every checkpoint to catch this is expensive. Measuring how the geometry of representations drifts on a fixed probe set is cheap and needs no labels, but it is unclear how fast that drift happens, what controls its speed, and whether it signals forgetting. We test four pre-registered claims: sub-1B models lose half their nearest-neighbour structure within 500 steps, isotropy declines monotonically, benchmark accuracy starts to fall once preservation drops below 60\%, and medical fine-tuning preserves more structure than code. We track 22 LoRA and full fine-tuning runs of five base models (160M--3B parameters, Qwen2.5 and Pythia) on medical, code, and general text at 17 checkpoints up to 3,000 steps. None of the four claims holds as stated. Larger models drift less at matched steps (medical, step 500: 0.42 for 160M \versus 0.85 for 3B), but the learning rate matters more than size: the same 0.5B model ends at 0.73, 0.57, or 0.22 preservation as the learning rate rises from 2e-4 to 1e-3. Every run that fell below 0.5 also had rising in-domain loss. No run becomes monotonically less isotropic. In Qwen models, benchmark forgetting begins while preservation is still 0.77--0.90. Code fine-tuning preserves more general-text geometry than medical fine-tuning (by 0.043, $p=3\times10^{-12}$), even though code shares less vocabulary with general text. Preservation tracks damage to general-text language modelling closely (Spearman $\rho=-0.90$), so it is a useful sign of unstable training. It does not provide a universal forgetting threshold.
